\subsection{NetBSD} \label{netbsd}

NetBSD is a \textbf{free, open-source, Unix-like operating system} founded in 1992 by Chris Demetriou, Theo de Raadt, Adam Glass,
and Charles M. Hannum \parencite{netBSDnd}. Its first public release, NetBSD 0.8, appeared in 1993 and was derived
from 4.3BSD Lite and 386BSD, both of which trace back to the Berkeley Software Distribution of Unix developed at UC Berkeley
\parencite{netBSDnd}. The BSD family as a whole has a well-documented lineage: NetBSD has its roots in both 386BSD
and 4.4BSD Lite, OpenBSD is based on the NetBSD source tree, FreeBSD 1.0 was based on 386BSD 0.1, and DragonFlyBSD began as
a fork of FreeBSD 4.8 \parencite{thomas_rantugi2022}. By 2022, the average age of these BSD operating systems was 26 years, with
NetBSD being among the oldest still under active development \parencite{thomas_rantugi2022}.

The hardware environment of the early 1990s directly shaped NetBSD's architectural decisions. At the time, there was no single
dominant processor architecture, memory was limited and expensive, and Unix systems were largely tied to specific vendor hardware.
These constraints pushed NetBSD's designers toward a highly abstracted kernel architecture. \textcite{chung2002} documents that the kernel's
I/O subsystem is built around the Virtual File System layer and the vnode interface, both of which are deliberately machine-independent
so that the system can run on a wide range of hardware without requiring deep rewrites. The memory management subsystem, known as UVM,
was designed to handle physical memory efficiently even on resource-constrained platforms, and the buffer cache was structured with
carefully managed free lists and hash tables to minimize memory overhead while keeping frequently accessed data available in memory \parencite{chung2002}.
The Berkeley Fast File System, which NetBSD uses as its primary filesystem, also reflects early hardware realities: it organizes disk storage into
cylinder groups to reduce mechanical seek times, which was a significant performance concern when spinning hard drives were the only available storage
medium \parencite{chung2002}.

The motivation for creating NetBSD came from dissatisfaction with the available options at the time. Commercial Unix systems were expensive
and locked to specific hardware vendors, while freely available alternatives like 386BSD existed but were not clean or stable enough
for widespread use. The founders wanted to build something that was both free and genuinely portable, meaning it could run on practically
any hardware platform without major rewrites \parencite{netBSDnd}. There was also a deliberate choice to ship NetBSD as a complete
operating system rather than just a kernel. This is an important distinction from Linux, which refers only to the kernel and does not include
userland utilities, filesystems, or windowing systems \parencite{thomas_rantugi2022}. By shipping a complete system, NetBSD became useful for both
research and embedded development, where having a single, coherent, well-documented codebase is more practical than assembling components from
multiple sources. \textcite{chung2002} notes that NetBSD's documentation quality and clean code structure made it a go-to reference for studying filesystem
implementation, particularly its FFS documentation in the 1.6 release.

NetBSD is primarily used in server, embedded systems, and research environments rather than general consumer desktop use. Its intended audience is
developers, embedded systems engineers, and researchers who need an operating system that can be ported to unusual or constrained hardware. Portability
is its defining priority, which sets it apart from OpenBSD, which prioritizes security and correctness, and FreeBSD, which leans toward
performance and enterprise features \parencite{thomas_rantugi2022}. This priority is reflected throughout the kernel: the machine-independent VFS
layer, the vnode interface, the platform-neutral memory management in UVM, and the Unified Buffer Cache that integrates the I/O and virtual
memory subsystems to avoid redundant data copies are all engineering choices made in service of portability and clean abstraction rather than
raw performance \parencite{chung2002}. NetBSD follows a monolithic kernel design, where the entire kernel runs in a single protected memory space. It
also adheres to the Unix philosophy of treating devices, sockets, and other system resources as files through a unified file descriptor interface,
which means application programs do not need to know whether they are reading from a disk, a network socket, or a hardware device since they interact
with all of them through the same system call layer \parencite{chung2002}.

In terms of influence, NetBSD draws from earlier BSD and Unix research, including the Sun Microsystems VFS design introduced in the
late 1980s and the Berkeley Fast File System, and it in turn directly influenced OpenBSD, which Theo de Raadt forked from the NetBSD
source tree following a disagreement with the core team \parencite{thomas_rantugi2022}. The broader BSD family also forms the foundation
of Apple's iOS, macOS, and watchOS \parencite{thomas_rantugi2022}, giving NetBSD's design decisions an indirect reach into consumer devices
worldwide. The development roles within NetBSD follow a clearly defined open-source model: contributors submit code, committers hold
write access to the repository, and a core team acts as the final gatekeepers of the source code \parencite{thomas_rantugi2022}. Code reviews
are not mandatory for all changes, but committers are expected to have nontrivial changes reviewed by at least one other person before
committing \parencite{thomas_rantugi2022}.

By 2022, NetBSD had 277 committers and 294,140 total commits, with a codebase of approximately 8.67 million lines of kernel code and
39.44 million lines of non-kernel code, the largest non-kernel codebase among the four BSD systems studied \parencite{thomas_rantugi2022}.
Its kernel code grows at a median annual rate of 7.15\% and non-kernel code at 6.19\%, with both subsystems increasing at a similar
pace \parencite{thomas_rantugi2022}. The distribution of commits also reveals something about its architecture: 52.68\% of commits touch
only kernel code, 46.04\% touch only non-kernel code, and just 1.27\% touch both simultaneously, the lowest mixed commit rate among
the four BSDs \parencite{thomas_rantugi2022}. This low percentage of mixed commits confirms that the kernel and userland are architecturally
well-separated, allowing engineers to work on each independently, which is exactly what a portability-first system requires. Kernel
commits in NetBSD are also consistently larger than non-kernel commits, with a median kernel commit size of 9 lines versus 7 lines for
non-kernel commits, and code reviews for kernel changes take longer to complete than non-kernel reviews, which is consistent with the
higher technical risk and complexity involved in modifying kernel code \parencite{thomas_rantugi2022}.
